\section{Incarnation of the Logos}

\begin{quotex}
In the beginning was the Word, and the Word was with God, and the Word was God. The same was in the beginning with God. All things were made by him: and without him was made nothing that was made. In him was life, and the life was the light of men. And the light shineth in darkness, and the darkness did not comprehend it. \flright{\textit{Gospel of John}}

\end{quotex}
In this gospel passage, the birth of Christ is expressed in theological language, since the mere historical fact has no meaning in itself. In his study on Hindu Doctrines, \textbf{Rene Guenon} devotes a chapter to the relationship between theology and metaphysics. In particular, he writes:

\begin{quotex}
Every theological truth, by means of a transposition dissociating it from its specific form, may be conceived in terms of the metaphysical truth corresponding to it, of which it is but a kind of translation…

\end{quotex}
This needs to be kept in mind, especially for those who are put off by theological language or can only see arbitrariness in theological disputes. Thus it is instructive to try to understand the metaphysical meaning of this passage, especially considering that it is the fundamental premise of Christianity. ``Word" is the translation of the Greek philosophical term Logos\footnote{\url{http://en.wikipedia.org/wiki/Logos}}, of which there is no really suitable translation; thus we shall leave it untranslated. There are two principles in play here.

\begin{itemize}
\item GOD: Principle of Being (Aquinas) 
\item LOGOS: Principle of order (Greek philosophy) 
\end{itemize}
The identity of God and the Logos thus implies that the principle of order is inseparable from Being itself. Were that not so, then chaos would be on equal footing with order, not the absence of order; ignorance would be equivalent to gnosis, not its opposite; revolution would be a possible option to hierarchy, not a revolt against being itself; darkness would be in perpetual opposition to light, rather than its shadow.



\flrightit{Posted on 2010-12-26 by Cologero }

\begin{center}* * *\end{center}

\begin{footnotesize}\begin{sffamily}



\texttt{Cologero on 2010-12-26 at 12:48 said: }

That is why the claims of Arius\footnote{\url{http://en.wikipedia.org/wiki/Arius}} that the Logos was created had to be opposed as heretical; this was not a hair-splitting theological game nor a power play.


\end{sffamily}\end{footnotesize}
